% ============================================================
%  Chapter 9 – Conclusion and Future Work
% ============================================================

% Relax line-breaking rules for this chapter so long \texttt{} identifiers
% (filenames, type names) can break gracefully across lines.
\sloppy

% ============================================================
\section{Summary of Achievements}
\label{sec:achievements}
% ============================================================

This thesis set out to transform a sequentially operated, perception-capable dual-arm
robotic cell into a fully parallelised, remotely monitored, and virtually validated
Industry~4.0 assembly system.  The following contributions were delivered and validated against measurable
criteria.

\subsection{Parallel execution and throughput}
A multithreaded asynchronous supervisor node, built on Python's \texttt{asyncio} library within
a ROS~2 \texttt{MultiThreadedExecutor}, was designed and deployed to orchestrate the ABB~IRB~120
and the AR4~Mk3 concurrently.  Hardware benchmarking demonstrated a \textbf{46.5\,\%
reduction in total assembly cycle time} (from 74.8\,s to 40.0\,s) and an \textbf{86.0\,\%
reduction in cumulative arm idle time} compared to the prior sequential architecture.

\subsection{Collision avoidance and zone management}
A dedicated \texttt{ZoneDetectionManager} node, operating at 500\,ms diagnostic intervals on a
separate callback group, continuously evaluated the Euclidean distance between end-effectors.
Across all forced stress tests the system maintained a minimum separation of \textbf{14.5\,cm}
— well above the 10\,cm safety threshold — and triggered MoveIt~Task~Constructor
(MTC) escape-route planning with a deterministic \textbf{average intervention latency of
0.089\,ms} (range: 0.075--0.114\,ms).

\subsection{Digital Twin: real-time monitoring}
A dual-mode Gazebo twin, launched via \texttt{twin.launch.py}, continuously received mirrored
joint-state commands from the physical robots through the \texttt{joint\_\allowbreak state\_\allowbreak to\_\allowbreak trajectory.py}
synchronisation nodes.  The divergence monitor maintained end-to-end visual lag below
\textbf{50\,ms}, enabling live spatial awareness without perceptible delay.  The joint divergence
safeguard issued system alerts at a 3.0\textdegree{} warning threshold and halted program
execution at a critical 8.0\textdegree{} anomaly limit, enforcing a continuous mathematical
audit of the physical--virtual consistency.

\subsection{Digital Twin: pre-execution validation}
The RViz teach panel and the \texttt{program\_\allowbreak sim\_\allowbreak runner.py} validation node together enforced
a strict \emph{validate-before-execute} contract.  Trajectory validation in the Gazebo twin
achieved a target precision profile of \textbf{RMSE\,$<$\,0.05\,rad} (2.86\textdegree{}) across
all tested waypoint sequences.  The \texttt{program\_\allowbreak real\_\allowbreak executor.py} node refused to command
the physical controllers unless the program-level \texttt{validated:~true} flag was present in
the YAML file, eliminating the possibility of sending an untested program to live hardware.

\subsection{High-speed software emergency-stop}
A dedicated e-stop channel (\texttt{/dt/estop}) allowed the remote operator to cancel active
trajectory goals with a \textbf{propagation and acknowledgement delay of $<$\,5\,ms},
providing deterministic remote failsafe capability during teleoperation.

\subsection{IIoT connectivity}
A two-channel communication architecture bridged two geographically separated laptops across
the public internet.  The \textbf{ZeroTier virtual overlay} provided NAT-traversing encrypted
connectivity; \textbf{Fast~DDS over unicast RTPS} carried real-time ROS~2 telemetry
with full QoS fidelity; and \textbf{Syncthing} delivered validated execution programs reliably
without any additional cloud infrastructure.

Table~\ref{tab:ch9_kpi_summary} consolidates the quantitative outcomes across all subsystems.

\begin{table}[H]
\centering
\caption{Key performance indicators achieved by the integrated system.}
\label{tab:ch9_kpi_summary}
\setlength{\tabcolsep}{5pt}
\begin{tabular}{lll}
\toprule
\textbf{Metric}                         & \textbf{Target}    & \textbf{Achieved}         \\
\midrule
Assembly cycle-time reduction           & $>$\,30\,\%        & 46.5\,\%                  \\
Arm idle-time reduction                 & $>$\,50\,\%        & 86.0\,\%                  \\
Minimum end-effector separation         & $\geq$\,10\,cm     & 14.5\,cm                  \\
Intervention trigger latency            & $<$\,1\,ms         & 0.089\,ms (avg)           \\
Digital twin visual lag                 & $<$\,100\,ms       & $<$\,50\,ms               \\
Pre-execution trajectory RMSE           & $<$\,0.10\,rad     & $<$\,0.05\,rad            \\
Software e-stop propagation delay       & $<$\,10\,ms        & $<$\,5\,ms                \\
\bottomrule
\end{tabular}
\end{table}

% ============================================================
\section{Challenges Overcome}
\label{sec:challenges}
% ============================================================

The path from individual-arm operation to a fully integrated Industry~4.0 cell involved
resolving a series of non-trivial technical obstacles, each of which shaped the final
architecture.

\subsection{Multithreading and ROS~2 callback isolation}
Running two independent arm controllers inside a single Python node required careful separation
of callback groups.  The initial design used a \texttt{SingleThreadedExecutor}, which caused
mutual blocking: a long-running MoveIt plan call for one arm starved the telemetry callbacks of
the other.  Migrating to a \texttt{MultiThreadedExecutor} with \texttt{ReentrantCallbackGroup}
for the arm action clients, and a \texttt{MutuallyExclusiveCallbackGroup} for the executive task
servers, eliminated contention while preserving the sequential semantics required within each
arm's own control flow.

\subsection{MoveIt~2 API compatibility under ROS~2 Jazzy}
The transition from ROS~2 Humble to Jazzy introduced header reorganisations in MoveIt~2,
deprecations of several \texttt{MoveGroupInterface} methods, and stricter requirements on
the \texttt{planning\_\allowbreak scene\_\allowbreak monitor} lifecycle.  Build failures arising from MOC (Qt
meta-object compiler) incompatibilities in the custom RViz~2 teach-panel plugin required the
CMake configuration to be restructured to generate MOC sources before the main library target.
Each API change was resolved incrementally against the Jazzy package manifests.

\subsection{Real-to-twin synchronisation fidelity}
Early monitoring-mode tests revealed that naïvely republishing real \texttt{/joint\_\allowbreak states}
directly to the Gazebo twin controllers produced discontinuous motion: the twin would jump
between positions rather than follow smoothly.  The solution was to route the measured joint
angles through \texttt{joint\_\allowbreak state\_\allowbreak to\_\allowbreak trajectory.py}, which constructs short
\texttt{FollowJointTrajectory} goals that the twin controllers execute at their native control
rate, producing fluid motion while faithfully reflecting the physical hardware.

\subsection{Cross-machine DDS discovery over the public internet}
ROS~2's default multicast-based RTPS discovery fails silently across NAT boundaries.  The
fix required authoring a custom Fast~DDS XML participant profile that disabled multicast,
declared explicit unicast initial-peer addresses (the ZeroTier virtual IPs), and whitelisted the
ZeroTier NIC exclusively.  An additional complication arose from the operating system
occasionally selecting the physical Wi-Fi interface for outbound UDP traffic even when the
ZeroTier NIC was specified; the whitelist resolved this by removing the ambiguity at the
middleware level.

\subsection{URDF collision conflicts in the twin description}
The combined \texttt{dual\_\allowbreak arms\_\allowbreak twin.xacro} initially produced a duplicate
\texttt{gripper\_\allowbreak base\_\allowbreak link} collision geometry arising from an overlap between the AR4 Mk3
gripper macro and the ABB gripper macro.  The conflict caused MoveIt's collision checker to
report spurious self-collisions at all configurations.  It was resolved by introducing
robot-specific prefix parameters into both macros and regenerating the SRDF allowed-collision
pairs.

\subsection{Sim-clock versus wall-clock in joint-value retrieval}
The digital twin environment enables Gazebo's simulation clock, which ROS~2 uses as the
source of \texttt{rclcpp::Clock::now()}.  The \texttt{getCurrentState()} call in MoveIt
uses this clock for timeout evaluation; with sim-clock paused or running at a non-unit rate,
the timeout fired immediately, causing IK calls to fail.  The fix required passing a
real-time clock instance explicitly to the MoveIt planning scene monitor, decoupling the
IK solver from the simulated time source.

% ============================================================
\section{Future Recommendations}
\label{sec:future}
% ============================================================

The system presented in this thesis lays a solid technical foundation upon which several
high-value extensions can be built.

\subsection{AI-driven path optimisation}
The current motion planner uses OMPL sampling-based algorithms (RRT-Connect) with fixed
timeout constraints.  Replacing or augmenting this with a learned motion-planning policy —
trained offline in the Gazebo twin using reinforcement learning — could reduce MTC escape-route
planning time for the AR4 from the measured \textasciitilde{}3\,s average to near-interactive
rates.  The twin's validated trajectory data provides a natural supervised-learning dataset.

\subsection{Force-aware handovers}
The handover process currently relies on position control and a fixed gripper-close delay.
Adding wrist force-torque sensors to both arms would enable compliant handover control: the
giver could release only after detecting the receiver's grip force exceed a threshold, and the
receiver could adapt its approach velocity based on the measured contact forces.  This would
eliminate the residual positional uncertainty at the handover point and extend the cell to
handle deformable or fragile workpieces.

\subsection{Advanced computer vision integration}
The current perception pipeline uses a single overhead RealSense~D435i together with an
ESP32-CAM for local IBVS.  Replacing the local camera with a calibrated wrist-mounted
depth sensor (e.g., Intel~RealSense~D405) on the AR4 would provide continuous 3D point-cloud
feedback during the final Cartesian approach, enabling real-time pose correction rather than
open-loop execution after the initial grasp pose is computed.

\subsection{Cloud-based analytics and predictive maintenance}
The IIoT architecture currently delivers telemetry to a single remote laptop.  Publishing
selected topics (joint torques, divergence values, cycle-time logs) to a cloud time-series
database (e.g., InfluxDB or Firebase Firestore) would enable long-term fleet-level analytics.
Anomaly detection models trained on historical divergence trends could flag impending mechanical
faults before they manifest as position errors, realising the predictive maintenance potential
discussed in the industrial feasibility analysis of Chapter~8.

\subsection{Formal safety certification pathway}
The current \texttt{ZoneDetectionManager} and software e-stop provide a strong engineering
safeguard, but they do not constitute a certified Safety-Rated Monitored Stop as defined by
ISO~10218 or ISO/TS~15066.  A natural next step is to integrate a certified safety PLC
(e.g., Pilz PNOZ~X or SICK~S3000) that independently monitors end-effector proximity via
laser safety scanners and enforces hard stops at the motor-drive level, making the cell
compliant with Speed and Separation Monitoring requirements for human proximity.

\subsection{Multi-robot scalability}
The current architecture supports exactly two heterogeneous arms.  Extending the
\texttt{UnifiedAssemblySupervisor} to manage $N>2$ arms — each with its own action server
and workspace zone definition — would require generalising the zone manager to an
$N$-body collision graph and the MTC handover pipeline to support multi-party transfers.
The layered DDS/ZeroTier communication stack scales transparently: additional participants
require only new entries in the \texttt{initialPeersList} of the Fast~DDS profile.

\vspace{1em}

In conclusion, this work demonstrates that a hybrid industrial and research-grade robotic cell
can be orchestrated as a coherent, parallelised, and remotely supervised Industry~4.0 system
using exclusively open-source middleware.  The digital twin layer — spanning real-time
monitoring, teach-panel programming, pre-execution validation, and divergence safeguards —
bridges the gap between simulation and physical execution in a manner that is both practically
useful and architecturally extensible.  The quantitative results confirm that the design choices
made at every layer, from the ZeroTier overlay to the YAML validation contract, translate
directly into measurable improvements in throughput, safety, and operational confidence.
